%% ============================================================
%% AULA 12 — Aplicações: interpolação e campos neurais
%% GBC073 — Inteligência Computacional (FACOM/UFU)
%%
%% Esqueleto com esboço de conteúdo — a desenvolver.
%% Ementa (ficha SEI 5116656): item 2 — aplicações: interpolação.
%%
%% Compilar:  latexmk -xelatex aula12.tex   (duas passadas!)
%% ============================================================
\documentclass[aspectratio=169, 11pt]{beamer}

\makeatletter
\def\input@path{{./}{../}}
\makeatother

\usetheme{InteligenciaComputacional}   % ANTES do babel: fontspec vem primeiro
\usepackage{babel}
\babelprovide[import, main]{brazilian}

\title[Aula 12 — Interpolação]{Aula 12: Aplicações — interpolação e campos neurais}
\subtitle[GBC073 — Inteligência Computacional]{Aproximação universal na prática \textbullet\ campos neurais \textbullet\ PINNs}
\author[Prof.\ Marcelo Albertini]{Prof.\ Marcelo Keese Albertini}
\institute{GBC073 — Inteligência Computacional \textbullet\ Universidade Federal de Uberlândia}
\date{2026}

\begin{document}

% ------------------------------------------------------------ capa
\begin{frame}[plain, noframenumbering]
  \titlepage
\end{frame}

% ------------------------------------------------------------ roteiro
\begin{frame}{Roteiro da aula}
  \begin{itemize}\setlength{\itemsep}{0.5ex}
    \item<1-> \textbf{Do teorema da aproximação universal à prática}
    \item<2-> \textbf{Campos neurais: coordenadas como entrada}
    \item<3-> \textbf{PINNs: a física como resíduo}
    \item<4-> \textbf{Quando a rede é o modelo do mundo}
  \end{itemize}
\end{frame}

% ------------------------------------------------------------
\section{Do teorema à prática}

\begin{frame}{Do teorema da aproximação universal à prática}
  \begin{itemize}\setlength{\itemsep}{0.4ex}
    \item Aula 2 provou: a rede \emph{pode} representar quase qualquer função.
    \item Interpolação é o teste mais direto: passar pelos pontos dados e
          generalizar entre eles.
    \item Redes interpolam suavemente, sem malha — e em dimensão alta.
    \item A ficha da disciplina pede exatamente isso: ``aplicações em
          interpolação''.
  \end{itemize}
\end{frame}

% ------------------------------------------------------------
\section{Campos neurais}

\begin{frame}{Campos neurais: coordenadas como entrada}
  \begin{itemize}\setlength{\itemsep}{0.4ex}
    \item Um campo neural (NeRF e família): a rede aprende uma função do
          espaço — a coordenada entra, a intensidade sai.
    \item Imagens, volumes 3D e cenas viram uma única rede contínua.
    \item O modelo é a memória: não se armazena o sinal, armazena-se a função.
    \item A coordenada sozinha, sem truque de codificação, subajusta — detalhe
          prático essencial.
  \end{itemize}
\end{frame}

% ------------------------------------------------------------
\section{PINNs}

\begin{frame}{PINNs: a física como resíduo}
  \begin{itemize}\setlength{\itemsep}{0.4ex}
    \item Redes informadas pela física (PINNs): a equação diferencial vira
          termo de perda.
    \item O resíduo da equação é avaliado em pontos amostrados do domínio.
    \item Sem malha, sem discretização — a rede é a solução.
    \item Fronteiras e condições iniciais entram como termos de perda
          adicionais.
  \end{itemize}
\end{frame}

% ------------------------------------------------------------
\section{A rede como modelo do mundo}

\begin{frame}{Quando a rede é o modelo do mundo}
  \begin{itemize}\setlength{\itemsep}{0.4ex}
    \item Do ajuste de curvas ao modelo de uma cena, de um material, de um
          sistema dinâmico.
    \item A rede como aproximador universal de \emph{processos}, não só de
          tabelas.
    \item Limites: extrapolação fora do domínio amostrado é aposta, não
          garantia.
  \end{itemize}
  \begin{quizbox}
    Um campo neural armazena o sinal discretizado na memória da GPU, como um
    arquivo de imagem.

    \smallskip
    \opcoes{Verdadeiro \;/\; Falso}
  \end{quizbox}
\end{frame}

\end{document}
